% =====================================================================
%  A GRASP algorithm for the Meal Delivery Routing Problem (MDRP)
%  Model specification extracted from:
%    D. Giraldo-Herrera and D. Alvarez-Martinez,
%    "A GRASP algorithm for the Meal Delivery Routing Problem", MIC 2024.
%  Department of Industrial Engineering, Universidad de Los Andes.
% =====================================================================
\documentclass[11pt]{article}

\usepackage[utf8]{inputenc}
\usepackage[T1]{fontenc}
\IfFileExists{lmodern.sty}{\usepackage{lmodern}}{}
\usepackage[margin=2.5cm]{geometry}
\usepackage{amsmath,amssymb}
\usepackage{booktabs}
\usepackage{algorithm}
\usepackage{algpseudocode}
\usepackage{enumitem}
\usepackage{hyperref}

% ---------------------------------------------------------------------
% Pseudocode helpers
% ---------------------------------------------------------------------
\algnewcommand\algorithmicinput{\textbf{Input:}}
\algnewcommand\Input{\item[\algorithmicinput]}
\algnewcommand\algorithmicparams{\textbf{Parameters:}}
\algnewcommand\Params{\item[\algorithmicparams]}
\newcommand{\var}[1]{\texttt{#1}}

\title{A GRASP Algorithm for the Meal Delivery Routing Problem\\[4pt]
       \large Model Specification}
\author{Daniel Giraldo-Herrera \and David \'Alvarez-Mart\'inez\\[2pt]
        \small Department of Industrial Engineering, Universidad de Los Andes, Bogot\'a, Colombia\\
        \small \texttt{ds.giraldoh@uniandes.edu.co}, \texttt{d.alvarezm@uniandes.edu.co}}
\date{}

\begin{document}
\maketitle

\begin{abstract}
This document contains the full specification of the GRASP metaheuristic
proposed for the Meal Delivery Routing Problem (MDRP): the problem notation,
the delivery assumptions, the three algorithmic components (constructive phase,
local search phase and the GRASP main loop) and the calibrated parameters.
\end{abstract}

% =====================================================================
\section{Problem description}
% =====================================================================

Users place orders throughout the day through a delivery application. Based on
the user's address and the active service radius, the platform decides whether
the user is inside the service area; if so, a delivery time is promised. The
platform dynamically adjusts the service radius according to expected and
observed demand. The objective is to determine feasible routes for couriers to
complete the pick-up and delivery of orders within time windows, optimizing one
or more performance measures (orders delivered, average delivery time, total
delivery time since the order was requested).

The MDRP is split into two parts:
\begin{enumerate}[label=(\roman*)]
  \item \textbf{Assignment:} assigning a courier to each order based on the
        geographical positions of the store, the courier and the customer.
  \item \textbf{Routing:} defining the route the courier follows to deliver
        orders in the shortest possible time, respecting the courier's maximum
        carrying capacity.
\end{enumerate}

% =====================================================================
\section{Notation}
% =====================================================================

\begin{table}[ht]
\centering
\begin{tabular}{@{}ll@{}}
\toprule
\textbf{Symbol} & \textbf{Description} \\
\midrule
\multicolumn{2}{@{}l}{\emph{Sets}} \\
$U$ & Set of users (customers) who place orders \\
$R$ & Set of restaurants (pick-up points) \\
$O$ & Set of orders placed by users \\
$C$ & Set of couriers \\
$S$ & Set of routes (bundles of orders) \\
\midrule
\multicolumn{2}{@{}l}{\emph{Users and restaurants}} \\
$\ell_u$ & Drop-off location of user $u \in U$ \\
$\ell_r$ & Pick-up location of restaurant $r \in R$ \\
\midrule
\multicolumn{2}{@{}l}{\emph{Orders} ($o \in O$)} \\
$r_o \in R$ & Restaurant associated with order $o$ \\
$u_o \in U$ & User associated with order $o$ \\
$\alpha_o$ & Placement time of order $o$ \\
$e_o$ & Ready time of order $o$ \\
$s^r$ & Pick-up service time (depends on the restaurant) \\
$s^u$ & Mandatory drop-off service time (linked to the user) \\
$c_o \in C$ & Courier assigned to order $o$ (an order may remain unassigned) \\
\midrule
\multicolumn{2}{@{}l}{\emph{Couriers} ($c \in C$)} \\
$e_c$ & On-time (shift start) of courier $c$ \\
$\ell_c$ & On-location (initial position) of courier $c$ \\
$l_c$ & Off-time (shift end) of courier $c$, with $l_c > e_c$ \\
$p_i$ & Compensation for each order (base amount + variable amount) \\
$m$ & Number of orders placed \\
\bottomrule
\end{tabular}
\caption{Notation of the MDRP.}
\label{tab:notation}
\end{table}

\paragraph{Courier compensation.} A courier's compensation at the end of the
shift is defined as
\begin{equation}
  \sum_{m=1}^{i} p_i \qquad \forall\, r \in R,
\end{equation}
where $p_i$ is the compensation for each order, composed of a base amount and a
variable amount that depends on the time it takes for the order to be accepted,
and $m$ is the number of orders placed.

\paragraph{Routes (bundles).} Orders from the same restaurant may be aggregated
into routes with a single pick-up and multiple drop-off locations. Every route
$s \in S$ must satisfy
\begin{equation}
  \bigl|\{\, r_o : o \in s \,\}\bigr| = 1,
\end{equation}
i.e., all orders in a route belong to the same restaurant. A courier carries at
most three orders per route. The information of the user associated with an
order is only disclosed to the courier once the courier accepts the order.

% =====================================================================
\section{Delivery assumptions}
% =====================================================================

\begin{enumerate}[label=A\arabic*.]
  \item Each order has a preparation time and a time window to be picked up and
        delivered; this time window cannot be exceeded.
  \item There may be unfulfilled orders due to time windows or to the absence
        of nearby couriers.
  \item If a courier arrives at the restaurant before the order is ready, the
        courier waits until the order is finished.
  \item A courier can carry at most 3 orders in a single route, provided they
        come from the same restaurant.
  \item Couriers use different vehicles with the following average velocities:
        \begin{center}
        \begin{tabular}{@{}lc@{}}
          \toprule
          \textbf{Vehicle} & \textbf{Average velocity (km/h)} \\
          \midrule
          Motorcycle & 20 \\
          Car        & 15 \\
          Bicycle    & 12 \\
          Walking    & 5  \\
          \bottomrule
        \end{tabular}
        \end{center}
  \item After delivering an order, the courier returns to his position and
        waits for another order.
  \item Distances are computed with the haversine formula and travel times
        with the average velocity of each courier's vehicle.
\end{enumerate}

% =====================================================================
\section{GRASP methodology}
% =====================================================================

GRASP (Greedy Randomized Adaptive Search Procedure), introduced by Feo and
Resende (1995), consists of two stages per iteration: a greedy randomized
\emph{construction} and a \emph{local search} improvement. Its behaviour is
governed by two parameters:
\begin{itemize}
  \item $\alpha \in [0,1]$: degree of randomness in the construction phase; it
        defines the size of the Restricted Candidate List (RCL) as a fraction
        of the sorted candidate list.
  \item \textit{iterations}: total number of constructed and improved
        solutions.
\end{itemize}

% ---------------------------------------------------------------------
\subsection{Constructive phase}
% ---------------------------------------------------------------------

The constructive phase iteratively assigns pending orders to available
couriers. For every courier--order pair it computes the travel time (distance
divided by the courier's velocity), sorts the candidates by travel time and
builds the RCL with the first $\alpha \cdot |\text{candidates}|$ elements. A
pair is selected at random from the RCL. If the selected courier has no orders
yet, the restaurant of the chosen order is recorded as the first element of the
courier's assignment. The process continues until no orders are pending.

\begin{algorithm}[H]
\caption{\textsc{Constructive\_Phase}}
\label{alg:construct}
\begin{algorithmic}[1]
\Input \var{couriers}, \var{orders}
\Params $\alpha$
\State $\var{assigned\_orders} \gets \emptyset$
\State $\var{couriers\_available} \gets \textsc{copy}(\var{couriers})$
\State $\var{orders\_pending} \gets \textsc{copy}(\var{orders})$
\While{$\var{orders\_pending} \neq \emptyset$}
  \State $\var{greedy\_values} \gets \emptyset$
  \For{$(\var{courier\_id}, \var{courier}) \in \var{couriers\_available}$}
    \For{$(\var{order\_id}, \var{order}) \in \var{orders\_pending}$}
      \State $\var{time\_travel} \gets
        \textsc{Calculate\_distance\_between}(\var{courier}, \var{order}) / \var{courier\_vel}$
      \State $\var{greedy\_values} \gets \var{greedy\_values} \cup
        \{(\var{time\_travel}, \var{courier\_id}, \var{order\_id})\}$
    \EndFor
  \EndFor
  \State sort \var{greedy\_values} by \var{time\_travel} (ascending)
  \State $\var{RCL} \gets$ first $\lceil \alpha \cdot |\var{greedy\_values}| \rceil$ elements of \var{greedy\_values}
  \State $(\var{selected\_courier\_id}, \var{selected\_order\_id}) \gets \textsc{random}(\var{RCL})$
  \If{$\var{selected\_courier\_id} \notin \var{assigned\_orders}$}
    \State $\var{current\_establishment} \gets \var{orders}[\var{selected\_order\_id}][\text{`restaurant\_ID'}]$
    \State $\var{assigned\_orders}[\var{selected\_courier\_id}] \gets [\var{current\_establishment}]$
  \EndIf
  \State $\var{assigned\_orders}[\var{selected\_courier\_id}] \gets \var{selected\_order\_id}$
  \State remove \var{selected\_order\_id} from \var{orders\_pending}
\EndWhile
\State \Return \var{assigned\_orders}
\end{algorithmic}
\end{algorithm}

% ---------------------------------------------------------------------
\subsection{Local search phase}
% ---------------------------------------------------------------------

The local search evaluates the total travel time of the assignment produced by
the constructive phase and explores \emph{swap} moves: for every pair of
distinct couriers, each order of the first courier is exchanged with each order
of the second. The index \texttt{[1:]} skips the first element of each list,
which stores the restaurant. If the new configuration reduces the total travel
time, it becomes the incumbent. The procedure ends when all swaps have been
explored.

\begin{algorithm}[H]
\caption{\textsc{Local\_Search\_Phase}}
\label{alg:ls}
\begin{algorithmic}[1]
\Input \var{couriers}, \var{orders}, \var{assigned\_orders}
\State $\var{better\_assign} \gets \textsc{copy}(\var{assigned\_orders})$
\State $\var{better\_time\_travel\_total} \gets \sum_{\text{assign} \in \var{assigned\_orders}} \var{time\_travel}$
\For{$(\var{courier\_id\_1}, \var{orders\_assigned\_1}) \in \var{assigned\_orders}$}
  \For{$(\var{courier\_id\_2}, \var{orders\_assigned\_2}) \in \var{assigned\_orders}$}
    \If{$\var{courier\_id\_1} \neq \var{courier\_id\_2}$}
      \For{$\var{order\_id\_1} \in \var{orders\_assigned\_1}[1{:}]$}
        \For{$\var{order\_id\_2} \in \var{orders\_assigned\_2}[1{:}]$}
          \State $\var{new\_assign} \gets \textsc{copy}(\var{assigned\_orders})$
          \If{$\var{order\_id\_1} \in \var{new\_assign}[\var{courier\_id\_1}]$}
            \State remove \var{order\_id\_1} from $\var{new\_assign}[\var{courier\_id\_1}]$
          \EndIf
          \If{$\var{order\_id\_2} \in \var{new\_assign}[\var{courier\_id\_2}]$}
            \State remove \var{order\_id\_2} from $\var{new\_assign}[\var{courier\_id\_2}]$
          \EndIf
          \State $\var{new\_assign}[\var{courier\_id\_1}] \gets [\var{order\_id\_2}]$
          \State $\var{new\_assign}[\var{courier\_id\_2}] \gets [\var{order\_id\_1}]$
          \State $\var{new\_time\_travel\_total} \gets \sum_{\text{assign} \in \var{new\_assign}} \var{time\_travel}$
          \If{$\var{new\_time\_travel\_total} < \var{better\_time\_travel\_total}$}
            \State $\var{better\_assign} \gets \textsc{copy}(\var{new\_assign})$
            \State $\var{better\_time\_travel\_total} \gets \var{new\_time\_travel\_total}$
          \EndIf
        \EndFor
      \EndFor
    \EndIf
  \EndFor
\EndFor
\State \Return \var{better\_assign}
\end{algorithmic}
\end{algorithm}

% ---------------------------------------------------------------------
\subsection{GRASP main loop}
% ---------------------------------------------------------------------

The main loop initializes an empty best solution with infinite cost and, for
the given number of iterations, runs the constructive phase followed by the
local search. The total travel time of the improved solution is compared with
the incumbent, which is updated whenever an improvement is found.

\begin{algorithm}[H]
\caption{\textsc{GRASP}}
\label{alg:grasp}
\begin{algorithmic}[1]
\Input \var{couriers}, \var{orders}
\Params $\alpha$, \var{iterations}
\State $\var{best\_solution} \gets \emptyset$
\State $\var{best\_cost} \gets \infty$
\For{$i = 1$ \textbf{to} \var{iterations}}
  \State $\var{constructed\_assign} \gets \textsc{Constructive\_Phase}(\var{couriers}, \var{orders};\, \alpha)$
  \State $\var{improved\_assign} \gets \textsc{Local\_Search\_Phase}(\var{couriers}, \var{orders}, \var{constructed\_assign})$
  \State $\var{current\_cost} \gets \sum_{\text{assign} \in \var{improved\_assign}} \var{time\_travel}$
  \If{$\var{current\_cost} < \var{best\_cost}$}
    \State $\var{best\_solution} \gets \var{improved\_assign}$
    \State $\var{best\_cost} \gets \var{current\_cost}$
  \EndIf
\EndFor
\State \Return \var{best\_solution}
\end{algorithmic}
\end{algorithm}

% =====================================================================
\section{Parameter calibration}
% =====================================================================

The platform runs its assignment and routing algorithms every two minutes, so
the GRASP must return a solution within that time. Calibration was performed as
follows:
\begin{itemize}
  \item $\alpha \in \{0, 0.1, 0.2, \dots, 1.0\}$ (step of $0.1$);
  \item $\textit{iterations} \in [500, 2000]$.
\end{itemize}
The best results were obtained with $\alpha \geq 0.7$ and close to 2000
iterations, but with computational times above five minutes. The selected
configuration is therefore
\begin{equation}
  \boxed{\alpha = 0.7, \qquad \textit{iterations} = 1000,}
\end{equation}
which yields solutions about 2\% below those obtained with more iterations, with
a computational time under two minutes.

\paragraph{Experimental setup.} Python implementation; Windows 11 (64-bit),
Intel Core i9-12900H CPU, 16\,GB RAM; 22 real-world one-day instances from a
Colombian delivery application. The GRASP is compared against the
simulation--optimization framework of Quintero Rojas (2020), using haversine
distances and average courier velocities in both approaches.

% =====================================================================
\section*{References}
% =====================================================================
\begin{enumerate}[label={[\arabic*]}, leftmargin=*]
  \item T.~A. Feo and M.~G.~C. Resende. Greedy randomized adaptive search
        procedures. \emph{Journal of Global Optimization}, 6:109--133, 1995.
  \item P.~Festa and M.~G.~C. Resende. GRASP. In R.~Mart\'i, P.~Pardalos and
        M.~Resende (eds.), \emph{Handbook of Heuristics}. Springer, Cham, 2016.
  \item S.~Quintero Rojas. \emph{Computational framework for solving the meal
        delivery routing problem}. Universidad de los Andes, 2020.
  \item D.~Reyes, A.~Erera, M.~Savelsbergh, S.~Sahasrabudhe and R.~O'Neil.
        The meal delivery routing problem. \emph{Optimization Online}, 6571,
        2018.
\end{enumerate}

\end{document}
